\documentclass[12pt]{article}
\usepackage[utf8]{inputenc}
\usepackage[T1]{fontenc}
\usepackage[margin=1in]{geometry}
\usepackage{setspace}
\usepackage{xcolor}
\usepackage[colorlinks=true, urlcolor=blue!60!black, linkcolor=blue!60!black]{hyperref}
\setlength{\parindent}{2em}
\setlength{\parskip}{1.25em}
\renewcommand{\baselinestretch}{1.0}
\begin{document}
\begin{center}
{\LARGE \textit{Michael H. v. Gerald D.}, 491 U.S. 110 (1989)\par}
\vspace{0.8em}
{\large Case Brief\par}
\end{center}
\vspace{1em}

\section{Memory Jogger}

In \href{https://www.courtlistener.com/opinion/112295/michael-h-v-gerald-d/}{\textit{Michael H. v. Gerald D.}}, a fractured Supreme Court upheld California's conclusive presumption that a child born to a married woman cohabiting with her husband is a child of the marriage, rejecting a biological father's due process claim and sparking a landmark debate over whether substantive due process rights must be defined at the most specific level of historical tradition.

\section{Detailed Case Facts}

On May 9, 1976, Carole D., an international model, married Gerald D., a top executive in a French oil company, in Las Vegas, Nevada, and the couple established a home in Playa del Rey, California (\href{https://supreme.justia.com/cases/federal/us/491/110/}{491 U.S. at 113}). In the summer of 1978, Carole began an adulterous affair with a neighbor, Michael H. In September 1980, Carole conceived a daughter, Victoria D., who was born on May 11, 1981 (\href{https://supreme.justia.com/cases/federal/us/491/110/}{491 U.S. at 113}). Gerald was listed as the father on Victoria's birth certificate and consistently held Victoria out to the world as his daughter, although Carole informed Michael shortly after the birth that she believed he might be the biological father (\href{https://supreme.justia.com/cases/federal/us/491/110/}{491 U.S. at 113--114}).

During the first three years of Victoria's life, she remained with Carole while living intermittently in several households (\href{https://supreme.justia.com/cases/federal/us/491/110/}{491 U.S. at 114}). In October 1981, Gerald moved to New York City for business while Carole stayed in California. Later that month, Carole, Michael, and Victoria underwent blood testing, which established a 98.07\% probability that Michael was Victoria's biological father (\href{https://supreme.justia.com/cases/federal/us/491/110/}{491 U.S. at 114}). In January 1982, Carole and Victoria visited Michael in St. Thomas, where Michael held Victoria out as his child. In March 1982, Carole left Michael and lived in California with another man, Scott K., before spending time with Gerald in New York and Europe during the spring and summer and returning to Scott in the fall (\href{https://supreme.justia.com/cases/federal/us/491/110/}{491 U.S. at 114}).

From August 1983 through May 1984, Carole again lived with Michael in Los Angeles whenever Michael was not in St. Thomas; Michael contributed to Victoria's support, held her out as his daughter, and Victoria called him ``Daddy'' (\href{https://supreme.justia.com/cases/federal/us/491/110/}{491 U.S. at 114, 143--144}). In April 1984, Carole and Michael signed a stipulation acknowledging Michael as Victoria's natural father, but Carole left Michael the next month and instructed her attorneys not to file the stipulation (\href{https://supreme.justia.com/cases/federal/us/491/110/}{491 U.S. at 114--115}). In June 1984, Carole reconciled permanently with Gerald and moved to New York, where they raised Victoria along with two additional children subsequently born into their marriage (\href{https://supreme.justia.com/cases/federal/us/491/110/}{491 U.S. at 115}).

The governing California statute, Cal. Evid. Code \S~621(a) (West Supp. 1989)---rooted in an 1872 California enactment---provided that ``the issue of a wife cohabiting with her husband, who is not impotent or sterile, is conclusively presumed to be a child of the marriage'' (\href{https://supreme.justia.com/cases/federal/us/491/110/}{491 U.S. at 115, 117}). Under \S~621(b)--(d), the presumption could be rebutted by blood tests only upon a motion filed within two years of the child's birth by the husband, or by the mother if the biological father filed an affidavit acknowledging paternity (\href{https://supreme.justia.com/cases/federal/us/491/110/}{491 U.S. at 115, 118}). Separately, Cal. Civ. Code \S~4601 authorized courts to award reasonable visitation rights to a parent unless detrimental to the child's best interests, and in the court's discretion to ``any other person having an interest in the welfare of the child'' (\href{https://supreme.justia.com/cases/federal/us/491/110/}{491 U.S. at 115--116, 133}).

\section{Procedural History}

In November 1982, after being rebuffed in his attempts to visit Victoria, Michael filed a filiation action in California Superior Court seeking to establish his paternity and visitation rights (\href{https://supreme.justia.com/cases/federal/us/491/110/}{491 U.S. at 114}). In March 1983, the court appointed an attorney and guardian ad litem for Victoria, who filed a cross-complaint asserting that if Victoria had more than one psychological or \textit{de facto} father, she was entitled to maintain filial relationships with both Michael and Gerald (\href{https://supreme.justia.com/cases/federal/us/491/110/}{491 U.S. at 114}). In May 1984, following a court-appointed psychologist's recommendation, the Superior Court granted Michael limited visitation privileges \textit{pendente lite} (\href{https://supreme.justia.com/cases/federal/us/491/110/}{491 U.S. at 115}).

On October 19, 1984, Gerald intervened and moved for summary judgment under Cal. Evid. Code \S~621. On January 28, 1985, the Superior Court granted summary judgment for Gerald after finding from Carole's and Gerald's affidavits that they were cohabiting at conception and birth and that Gerald was neither impotent nor sterile (\href{https://supreme.justia.com/cases/federal/us/491/110/}{491 U.S. at 115}). The Superior Court rejected Michael's and Victoria's constitutional challenges and denied motions for continued visitation pending appeal under Cal. Civ. Code \S~4601 (\href{https://supreme.justia.com/cases/federal/us/491/110/}{491 U.S. at 115--116}).

The California Court of Appeal, Second Appellate District, affirmed (\textit{Michael H. v. Gerald D.}, 191 Cal. App. 3d 995, 236 Cal. Rptr. 810 (1987)), rejecting Michael's procedural and substantive due process claims, Victoria's due process and equal protection claims, and Victoria's claim to visitation under \S~4601 (\href{https://supreme.justia.com/cases/federal/us/491/110/}{491 U.S. at 116}). On July 30, 1987, the California Supreme Court denied discretionary review (\href{https://supreme.justia.com/cases/federal/us/491/110/}{491 U.S. at 116}). The U.S. Supreme Court noted probable jurisdiction on February 29, 1988 (485 U.S. 903), heard oral argument on October 11, 1988, and issued its decision on June 15, 1989 (\href{https://tile.loc.gov/storage-services/service/ll/usrep/usrep491/usrep491110/usrep491110.pdf}{491 U.S. at 110, 116}). The Court declined to reach Michael's equal protection claim because it was neither raised nor passed upon below (\href{https://supreme.justia.com/cases/federal/us/491/110/}{491 U.S. at 116--117}).

\section{Judicial Votes}

The Supreme Court affirmed the judgment by a 5--4 vote without a majority opinion (\href{https://supreme.justia.com/cases/federal/us/491/110/}{491 U.S. at 112--113}):
\begin{itemize}
\item \textbf{Plurality Opinion (Scalia, J.):} Announced the judgment of the Court and delivered an opinion joined in full by \textbf{Rehnquist, C.J.}, and joined in all except Footnote 6 by \textbf{O'Connor, J.} and \textbf{Kennedy, J.} (4 votes for the plurality reasoning; only 2 votes for Footnote 6).
\item \textbf{Concurring in Part (O'Connor, J., joined by Kennedy, J.):} Joined Justice Scalia's opinion except for Footnote 6, rejecting Footnote 6's rigid rule requiring historical traditions to be defined at ``the most specific level'' available (\href{https://supreme.justia.com/cases/federal/us/491/110/}{491 U.S. at 132}).
\item \textbf{Concurring in the Judgment (Stevens, J.):} Provided the fifth vote to affirm, assuming Michael might have a constitutionally protected liberty interest but concluding that Cal. Civ. Code \S~4601 gave him a constitutionally adequate opportunity to seek visitation as an ``other person'' (\href{https://supreme.justia.com/cases/federal/us/491/110/}{491 U.S. at 132--136}).
\item \textbf{Dissenting Opinion (Brennan, J., joined by Marshall, J., and Blackmun, J.):} Argued that Michael possessed a protected liberty interest under the unwed-father precedents and that \S~621 violated procedural due process by terminating that interest without a hearing (\href{https://supreme.justia.com/cases/federal/us/491/110/}{491 U.S. at 136--157}).
\item \textbf{Dissenting Opinion (White, J., joined by Brennan, J.):} Argued that biological paternity combined with a developed parental relationship created a protected liberty interest and that \S~621 denied Michael a meaningful opportunity to be heard (\href{https://supreme.justia.com/cases/federal/us/491/110/}{491 U.S. at 157--163}).
\end{itemize}

\section{Holding}

\textbf{Judgment:} The judgment of the California Court of Appeal affirming summary judgment for Gerald D. and upholding the application of Cal. Evid. Code \S~621 was \textbf{affirmed} (5--4) (\href{https://supreme.justia.com/cases/federal/us/491/110/}{491 U.S. at 110, 132}).

\textbf{Rule Announced:} Because there was no majority opinion, the narrowest ground supporting the judgment is that a state does not violate the Due Process Clause of the Fourteenth Amendment by applying a conclusive statutory presumption of marital paternity (Cal. Evid. Code \S~621) to bar a biological father from establishing legal paternity over a child conceived within and born into an extant, cohabiting marriage where both spouses wish to raise the child as their own---at least where state law does not deprive the natural father of a constitutionally fair process or where no fundamental right to override the marital family exists (\href{https://supreme.justia.com/cases/federal/us/491/110/}{491 U.S. at 129--130, 132--136}). The plurality also concluded that the child has neither a due process right to dual fatherhood nor an equal protection right to rebut the presumption under rational basis review (\href{https://supreme.justia.com/cases/federal/us/491/110/}{491 U.S. at 130--132}).

\textbf{Holding vs. Dicta:} The \textbf{holding} is limited to upholding California's statutory scheme as applied to a child born into an extant marital union that wishes to embrace the child. Crucially, Justice Scalia's famous \textbf{Footnote 6} methodology---requiring courts in substantive due process cases to define asserted rights at ``the most specific level at which a relevant tradition protecting, or denying protection to, the asserted right can be identified''---commanded only \textbf{two votes} (Scalia and Rehnquist) and is therefore non-binding \textbf{plurality/minority dicta}, expressly rejected by seven Justices (O'Connor, Kennedy, Stevens, Brennan, Marshall, Blackmun, and White) (\href{https://supreme.justia.com/cases/federal/us/491/110/}{491 U.S. at 127--128 n.6, 132, 136--137}). Similarly, the plurality's Footnote 7 reservation regarding situations where the marital parents do \textit{not} wish to raise the child jointly is expressly non-decisional dicta (\href{https://supreme.justia.com/cases/federal/us/491/110/}{491 U.S. at 129 n.7}).

\section{Analysis of Opinions}

\subsection{Plurality Opinion (Justice Scalia, joined by Chief Justice Rehnquist, and in all but Footnote 6 by Justices O'Connor and Kennedy)}

Justice Scalia first rejected Michael's procedural due process challenge, explaining that although Cal. Evid. Code \S~621 is phrased as an ``irrebuttable presumption,'' it actually implements a \textbf{substantive rule of law}: California has decided as a matter of social policy that biological paternity by an outsider is legally irrelevant when a child is conceived during and born into an existing cohabiting marriage (\href{https://supreme.justia.com/cases/federal/us/491/110/}{491 U.S. at 119--121}). Calling a substantive classification an irrebuttable presumption does not convert a substantive challenge into a procedural one; under cases like \textit{Weinberger v. Salfi}, 422 U.S. 749 (1975), the question is the substantive validity of the classification itself (\href{https://supreme.justia.com/cases/federal/us/491/110/}{491 U.S. at 120--121}).

Turning to substantive due process, Justice Scalia warned that extending ``liberty'' beyond freedom from physical restraint risks turning judicial review into judges' personal predilections (\href{https://supreme.justia.com/cases/federal/us/491/110/}{491 U.S. at 121--122}). To constrain judicial discretion, an asserted liberty interest must be not only fundamental but ``traditionally protected by our society'' (\href{https://supreme.justia.com/cases/federal/us/491/110/}{491 U.S. at 122}). Scalia reinterpreted the unwed-father precedents---\textit{Stanley v. Illinois}, 405 U.S. 645 (1972), \textit{Quilloin v. Walcott}, 434 U.S. 246 (1978), \textit{Caban v. Mohammed}, 441 U.S. 380 (1979), and \textit{Lehr v. Robertson}, 463 U.S. 248 (1983)---as resting not on biology plus a developed relationship in isolation, but on historic respect for the ``unitary family'' (\href{https://supreme.justia.com/cases/federal/us/491/110/}{491 U.S. at 123}). Surveying common law treatises (Bracton, Blackstone, Kent, and Schouler) and modern state statutes, Scalia found that Anglo-American tradition affirmatively protected the marital family against paternity claims by an ``adulterous natural father'' (\href{https://supreme.justia.com/cases/federal/us/491/110/}{491 U.S. at 124--127}).

In \textbf{Footnote 6} (joined only by Chief Justice Rehnquist), Scalia defended his historical method against Justice Brennan's accusation of arbitrariness, asserting that courts must consult ``the most specific level at which a relevant tradition protecting, or denying protection to, the asserted right can be identified'' (\href{https://supreme.justia.com/cases/federal/us/491/110/}{491 U.S. at 127--128 n.6}). Finally, in Part IV, the plurality rejected Victoria's due process claim to ``dual fatherhood'' as unsupported by history and rejected her equal protection claim under rational basis review because she was treated as a legitimate child and limiting rebuttal to spouses rationally protects marital peace (\href{https://supreme.justia.com/cases/federal/us/491/110/}{491 U.S. at 130--132}).

\subsection{Concurrence in Part (Justice O'Connor, joined by Justice Kennedy)}

Justice O'Connor joined Justice Scalia's opinion in full \textit{except} Footnote 6. In a brief but consequential concurrence, she cautioned that Footnote 6's requirement of a single ``most specific level'' historical test ``may be somewhat inconsistent with our past decisions in this area,'' citing \textit{Griswold v. Connecticut}, 381 U.S. 479 (1965), \textit{Eisenstadt v. Baird}, 405 U.S. 438 (1972), \textit{Loving v. Virginia}, 388 U.S. 1 (1967), and \textit{Turner v. Safley}, 482 U.S. 78 (1987) (\href{https://supreme.justia.com/cases/federal/us/491/110/}{491 U.S. at 132}). Citing Justice Harlan's dissent in \textit{Poe v. Ullman}, 367 U.S. 497 (1961), O'Connor refused to ``foreclose the unanticipated by the prior imposition of a single mode of historical analysis'' (\href{https://supreme.justia.com/cases/federal/us/491/110/}{491 U.S. at 132}).

\subsection{Concurrence in the Judgment (Justice Stevens)}

Justice Stevens agreed with the plurality that the Constitution does not require a state to declare biological facts when no legal consequence follows, but he rejected Scalia's categorical conclusion that a natural father can never have a constitutionally protected liberty interest in a child born into another man's marriage (\href{https://supreme.justia.com/cases/federal/us/491/110/}{491 U.S. at 132--133}). Assuming \textit{arguendo} that Michael's relationship with Victoria was strong enough to confer a constitutional right to seek visitation, Stevens concluded that Cal. Civ. Code \S~4601---which authorizes discretionary visitation for ``any other person having an interest in the welfare of the child''---provided Michael a fair opportunity to be heard, and that the trial judge did not act fundamentally unfairly in deferring to the mother's judgment given the restored stability of Carole and Gerald's marriage (\href{https://supreme.justia.com/cases/federal/us/491/110/}{491 U.S. at 133--136}).

\subsection{Dissenting Opinion (Justice Brennan, joined by Justices Marshall and Blackmun)}

Justice Brennan opened by counting noses: five Justices (Stevens, White, Brennan, Marshall, and Blackmun) refused to foreclose a natural father's liberty interest in this setting, and only two Justices endorsed Footnote 6 (\href{https://supreme.justia.com/cases/federal/us/491/110/}{491 U.S. at 136--137}). Brennan attacked Scalia's reliance on specific traditions as both subjective and contrary to precedent, noting that \textit{Griswold}, \textit{Eisenstadt}, \textit{Ingraham v. Wright}, and \textit{Stanley} would all have come out differently under Footnote 6 (\href{https://supreme.justia.com/cases/federal/us/491/110/}{491 U.S. at 137--141}). Restricting the Due Process Clause to practices already protected by majorities turns the Fourteenth Amendment ``into a redundancy'' and transforms a living Constitution into a ``stagnant, archaic, hidebound document'' (\href{https://supreme.justia.com/cases/federal/us/491/110/}{491 U.S. at 140--141}).

On the merits, Brennan argued that \textit{Stanley}, \textit{Quilloin}, \textit{Caban}, and \textit{Lehr} establish a clear principle: biological paternity plus a substantial developed parent-child relationship creates a protected liberty interest (\href{https://supreme.justia.com/cases/federal/us/491/110/}{491 U.S. at 142--143}). By defining ``unitary family'' to exclude Michael, Carole, and Victoria, the plurality conflated the threshold existence of a liberty interest with the State's countervailing interest in terminating it (\href{https://supreme.justia.com/cases/federal/us/491/110/}{491 U.S. at 145--147}). Brennan also dismissed Stevens's reading of \S~4601 as ``mere wishful thinking'' under California precedent (\textit{Vincent B. v. Joan R.}), concluding that \S~621 unconstitutionally denied Michael a procedural due process hearing to prove paternity (\href{https://supreme.justia.com/cases/federal/us/491/110/}{491 U.S. at 148--157}).

\subsection{Dissenting Opinion (Justice White, joined by Justice Brennan)}

Justice White focused squarely on the unwed-father precedents and procedural due process. He argued that Michael clearly satisfied \textit{Lehr v. Robertson}'s demanding standard: unlike Mr. Lehr, Michael lived with Victoria, held her out as his child, supported her financially, and was called ``Daddy'' (\href{https://supreme.justia.com/cases/federal/us/491/110/}{491 U.S. at 158--160}). White rejected California's justifications for \S~621---noting that blood tests eliminate ancient evidentiary uncertainties, that stigma of illegitimacy is weak where a father seeks to claim his child, and that Gerald already knew of the affair (\href{https://supreme.justia.com/cases/federal/us/491/110/}{491 U.S. at 161--162}). Because \S~621 barred Michael at the outset from introducing blood-test evidence and rendered him ``a stranger to his child'' without a meaningful opportunity to be heard, White concluded it violated the Fourteenth Amendment (\href{https://supreme.justia.com/cases/federal/us/491/110/}{491 U.S. at 162--163}).

\section{Examples: Future Applications}

\begin{enumerate}
\item \textbf{Same-Side Application (Upholding Marital Presumption Against Outside Biological Father):} A married couple lives together continuously throughout pregnancy and the child's early years, and both spouses wish to raise the child as their own. Three years after the birth, a former coworker with whom the wife had a brief affair files suit seeking DNA testing and parental custody over the couple's objection. Under \textit{Michael H.}, a state statute conclusively presuming the husband to be the legal father is constitutional because the child was conceived within and born into an extant, cohabiting marital family that wishes to embrace the child.

\item \textbf{Same-Side Application (Rejecting a Child's Claim to Legal Recognition of Three Parents):} A child born into an intact marriage sues through a guardian ad litem claiming a Fourteenth Amendment substantive due process right to maintain simultaneous legal parentage with her mother, her mother's husband, and her biological father. Under Part IV of \textit{Michael H.}, the claim fails because Anglo-American legal traditions do not mandate state recognition of multiple legal fathers, and limiting standing to rebut marital legitimacy to the spouses survives rational basis equal protection review.

\item \textbf{Opposite-Side Application (Unwed Biological Father Where Mother Is Unmarried):} An unmarried couple lives together for six years and raises two children together; the father financially supports the children and participates daily in their care. After the mother dies, the state automatically places the children in foster care under a statute conclusively presuming unwed fathers unfit without a hearing. Because there is no competing intact marital family, \textit{Michael H.} does not apply; instead, \textit{Stanley v. Illinois} governs, and the biological father possesses a constitutionally protected liberty interest requiring an individualized fitness hearing.

\item \textbf{Opposite-Side Application (State Statute Depriving Biological Father of Parental Rights When Marriage Dissolved Before Birth):} A wife separates permanently from her husband, files for divorce, lives exclusively with the biological father for two years before and after the child's birth, and the husband expressly disclaims paternity and refuses to raise the child. If the state nevertheless bars the biological father from establishing paternity solely because the divorce was not finalized until after conception, \textit{Michael H.} does not control: Scalia's plurality expressly limited its holding in Footnote 7 to cases where the husband and wife cohabited and jointly wish to raise the child, and five Justices in \textit{Michael H.} refused to foreclose a natural father's liberty interest in non-intact marital settings.

\item \textbf{Fence-Sitter Application (Known Sperm Donor or Biological Father Seeking Visitation Against a Married Same-Sex Couple):} A married couple conceives a child using a known donor without formal clinical intermediation; the donor participates regularly in the child's first two years of life with the couple's consent, after which the married couple cuts off contact and invokes the state's marital parentage presumption. This is a fence-sitter: on one hand, \textit{Michael H.} supports a state's categorical preference for the intact marital family unit over an outside biological progenitor; on the other hand, if the married parents intentionally invited the donor into a co-parenting role and the state provides no discretionary ``other person'' visitation avenue like Cal. Civ. Code \S~4601 (which proved decisive for Justice Stevens's fifth vote), a court could find a due process violation or distinguish \textit{Michael H.}'s ``adulterous natural father'' framing.
\end{enumerate}

\section{Critique}

\textit{Michael H. v. Gerald D.} sits at the epicenter of modern constitutional theory regarding unenumerated rights and the ``level of generality'' problem. From an \textbf{originalist and textualist perspective}, defenders of Justice Scalia's approach argue that because the text of the Fourteenth Amendment's Due Process Clause guarantees ``process'' rather than open-ended substantive liberties, any judicial recognition of unenumerated fundamental rights must be tightly disciplined by concrete historical practice (\href{https://texaslawreview.org/in-defense-of-substantive-due-process/}{Texas Law Review}; \href{https://journals.law.harvard.edu/jlpp/wp-content/uploads/sites/90/2026/03/57_14HarvJLPubPoly8531991.pdf}{Harvard Journal of Law \& Public Policy}). Scholars analyzing traditionalism note that consulting the most specific identifiable tradition seeks to prevent judges from smuggling personal moral preferences into constitutional law under abstract labels like ``family'' or ``autonomy'' (\href{https://repository.law.umich.edu/cgi/viewcontent.cgi?article=1371&context=mlr}{Michigan Law Review}; \href{https://scholarship.law.gwu.edu/cgi/viewcontent.cgi?article=2519&context=faculty_publications}{GWU Law Scholarly Commons}).

From a \textbf{progressive and doctrinal perspective}, however, Footnote 6 drew immediate and sustained scholarly fire. In their landmark critique, Laurence H. Tribe and Michael C. Dorf demonstrated that selecting a ``level of generality'' is never value-neutral: even defining a tradition at its ``most specific level'' requires normative choices about which historical facts are relevant (for example, adultery versus cohabitation, or blood testing versus common-law evidentiary limits) (\href{https://scholarship.law.cornell.edu/cgi/viewcontent.cgi?article=1121&context=facpub}{Tribe \& Dorf, Levels of Generality in the Definition of Rights}). Critics also emphasize that tradition-based substantive due process fails to genuinely constrain judicial discretion because judges retain wide latitude over whose traditions count and when a tradition becomes obsolete (\href{https://www.law.nyu.edu/sites/default/files/ECM_PRO_062707.pdf}{NYU Law}; \href{https://scholarship.law.duke.edu/cgi/viewcontent.cgi?article=3364&context=dlj}{Duke Law Journal}). Family-law and children's-rights scholars further criticize the plurality for elevating a rigid, state-approved marital form over the lived psychological attachments of both the child and a committed biological parent (\href{https://scholarworks.law.ubalt.edu/cgi/viewcontent.cgi?article=2129&context=all_fac}{University of Baltimore Law}).

Ultimately, \textit{Michael H.} exposes a deep clash of \textbf{underlying competing values}: judicial restraint and democratic majoritarianism (leaving contentious family-law trade-offs to state legislatures) versus counter-majoritarian protection of individual liberty and nonconforming family bonds; and institutional protection of the stable marital household versus individualized, fact-sensitive fairness for biological parents and children whose lived realities defy traditional categories.

\section{Key Quotations}

\begin{quote}
``As we view them, they rest not upon such isolated factors but upon the historic respect --- indeed, sanctity would not be too strong a term --- traditionally accorded to the relationships that develop within the unitary family.''
\end{quote}
\noindent --- Justice Scalia, plurality opinion, \href{https://supreme.justia.com/cases/federal/us/491/110/}{491 U.S. at 123} (\href{https://www.courtlistener.com/opinion/112295/michael-h-v-gerald-d/}{CourtListener}).

\begin{quote}
``We refer to the most specific level at which a relevant tradition protecting, or denying protection to, the asserted right can be identified. ... Although assuredly having the virtue (if it be that) of leaving judges free to decide as they think best when the unanticipated occurs, a rule of law that binds neither by text nor by any particular, identifiable tradition is no rule of law at all.''
\end{quote}
\noindent --- Justice Scalia, plurality opinion (joined on Footnote 6 only by Chief Justice Rehnquist), \href{https://supreme.justia.com/cases/federal/us/491/110/}{491 U.S. at 127--128 n.6} (\href{https://tile.loc.gov/storage-services/service/ll/usrep/usrep491/usrep491110/usrep491110.pdf}{U.S. Reports}).

\begin{quote}
``On occasion, the Court has characterized relevant traditions protecting asserted rights at levels of generality that might not be `the most specific level' available. ... I would not foreclose the unanticipated by the prior imposition of a single mode of historical analysis.''
\end{quote}
\noindent --- Justice O'Connor, concurring in part (joined by Justice Kennedy), \href{https://supreme.justia.com/cases/federal/us/491/110/}{491 U.S. at 132} (\href{https://www.courtlistener.com/opinion/112295/michael-h-v-gerald-d/}{CourtListener}).

\begin{quote}
``The document that the plurality construes today is unfamiliar to me. It is not the living charter that I have taken to be our Constitution; it is instead a stagnant, archaic, hidebound document steeped in the prejudices and superstitions of a time long past.''
\end{quote}
\noindent --- Justice Brennan, dissenting (joined by Justices Marshall and Blackmun), \href{https://supreme.justia.com/cases/federal/us/491/110/}{491 U.S. at 141} (\href{https://www.courtlistener.com/opinion/112295/michael-h-v-gerald-d/}{CourtListener}).

\begin{quote}
``Michael eagerly grasped the opportunity to have a relationship with his daughter (he lived with her; he declared her to be his child; he provided financial support for her) and still, with today's opinion, his opportunity has vanished. He has been rendered a stranger to his child.''
\end{quote}
\noindent --- Justice White, dissenting (joined by Justice Brennan), \href{https://supreme.justia.com/cases/federal/us/491/110/}{491 U.S. at 163} (\href{https://www.courtlistener.com/opinion/112295/michael-h-v-gerald-d/}{CourtListener}).

\end{document}
